\documentclass[10pt,letterpaper]{article}
\usepackage{cornell-notes}

\title{Introduction}
\author{}
\date{}
\setDocTitle{Introduction}
\setDocSubtitle{C++ 2024}
\setDocOwner{C++ 2024 Cornell Notes}
\setDocDate{\today}
\setCornellCollection{C++ 2024 Cornell Notes}
\setCornellUnitType{Introduction}
\setCornellUnitNumber{Introduction}
\setCornellUnitTitle{Introduction}
\hypersetup{pdftitle={Introduction},pdfsubject={Cornell notes},pdfauthor={}}

\begin{document}
\maketitle
\begin{CornellOverviewBox}{Overview}
\textbf{Classification:} Introductory\\[2pt]
\textbf{Learning objective:} Understand the rules, terminology, and semantic consequences associated with Introduction.
\end{CornellOverviewBox}\begin{CornellSummaryBox}{Summary}
{\sffamily\bfseries\color{Accent} Summary}\par\vspace{3pt}The introduction explains how the technical material is organized and referenced. The section supplies the rules and semantic distinctions needed to interpret this topic in context with its cross-references. Apply the section together with its cited rules so that terminology and constraints are interpreted in context.
\end{CornellSummaryBox}

\end{document}
